\todo{note: this is cut content}

\begin{comment}
The simplest method for comparing work done via different algorithms is for all reflecting chains to measure that work via a common unit. How can we do this? We could try to measure it in some external unit like USD. However, that would require accounting for external factors like the availability of silicon and mining rigs, the cost of electricity, and various exchange rates. Accounting for those external factors impractical, so whatever method we choose, those factors must *cancel out*.
\end{comment}

In blockchains like Bitcoin and Ethereum, block rewards are denominated
in the \emph{root token} of that chain. What if we had two chains with
the same root token?

Consider two PoW blockchains that share a root token, but use different
hashing algorithms. We can't directly compare their rates of work
because the PoW algorithms are different, but we \emph{can} compare
their rates of \emph{token generation via block rewards} (normalized
against time). Why does this work? Market externalities are encapsulated
by the \emph{economic} decisions that miners have to make. Thus, the
rate of work of each chain will approach equilibrium, and if both are at
equilibrium we can say that the work contributed to each chain is in
proportion to the token generation rate. What do I mean by this? Let's
look at an example.

Say we had Chain L and Chain R: two PoW chains using different
algorithms, but with the \emph{same} root token and the same block
production rate. Firstly, we need a way to determine what their block
rewards are, and the intuitive solution is to set each chain's block
reward proportionally to the percentage of root tokens (i.e.,
\emph{coins}) on that chain. Say that 100\% of the root tokens should
correspond to 50 coins generated per block. So, if one chain had 60\% of
the root tokens, then 30 coins are generated per block \emph{on that
chain}, and (since we only have 2 chains) \emph{the other chain} has 20
coins generated per block --- corresponding to the remaining 40\% of root
tokens. One reason this method makes sense is that the \emph{ratio} of
coins on each chain is not affected by block rewards. We can generalize
the method by normalizing against time so that we're comparing the rates
of coin production rather than the block rewards themselves.

Now that we know what the block rewards are, and have defined them in
terms of the percentage of total coins that are on that chain, we can
work on comparing the chains' hash-rates. What sort of foundation could
we do this from? What about \emph{equal work for equal reward}? Because
we have defined block rewards in terms of \emph{where} root tokens are
held, we can measure things like \emph{hashes per token} (when
considering block rewards particularly). Crucially, we can measure this
\emph{for each chain}, which allows us to --- \emph{contextually} --- make
claims like 100 hashes on Chain L are worth 25 hashes on Chain R.

Since we know the percentage of root tokens on each chain for each
moment in history, we can safely use that figure in chain-weight
calculations. The reliability of that data will be the same as the
reliability of the blockchains themselves, provided we enforce the 2-way
peg that ensures no root tokens are created or destroyed (which would
violate protocol rules).

\autoref{alg:weightof-1} and \autoref{alg:weightof-ratio} detail
simplistic \textsc{WeightOf} functions that return a weight normalized
in coins (i.e., root tokens). It does not account for some things, e.g.,
changes to the difficulty of \(C\) over time.

\begin{algorithm}
\caption{A simplistic \textsc{WeightOf} for networks of a single root token with equal block frequencies}\label{alg:weightof-1}
\begin{algorithmic}
\Procedure{WeightOf}{$B_i, state$} \Comment{The weight of a block normalized to coins}
    \State $t \gets$ \Call{BlockTimestamp}{$B_i$}
    \State $C \gets$ \Call{ChainOf}{$B_i$, $state$} \Comment{Either the local or a reflecting chain}
    \State $C_r \gets$ \Call{BlockRewardOfChainAt}{$C$, $t$, $state$} \Comment{Block reward of $C$ at $t$}
    \State \Return{$C_r$}
\EndProcedure
\end{algorithmic}
\end{algorithm}


\IfPackageLoadedTF{lwarp}{}{
\begin{algorithm}
\caption{A \textsc{WeightOf} function for networks of a single root token based on the ratio of root tokens that it hosts}
\label{alg:weightof-ratio}
\begin{algorithmic}
\Procedure{WeightOf}{$B_i, state$} \Comment{Returns L-hashes}
    \State $t \gets$ \Call{BlockTimestamp}{$B_i$}
    \State $C \gets$ \Call{ChainOf}{$B_i$, $state$}
    \State $C_t \gets$ \Call{RootTokensOnChainAt}{$C$, $t$, $state$} \Comment{Number of root tokenson $C$ at $t$}
    \State $C_f \gets$ \Call{BlockFrequencyOfChain}{$C$, $state$} \Comment{Block production Hz of $C$}
    \State $L_t \gets$ \Call{LocalRootTokens}{$t$, $state$}
    \State $L_f \gets$ \Call{LocalBlockBlockFrequency}{$t$, $state$}
    \State $L_d \gets$ \Call{LocalBlockDifficultyAt}{$t$, $state$}
    \State \Return{$R_t \cdot L_f \cdot L_d \div L_t \div R_f$} \Comment{See \autoref{eq:conv-srt}}
\EndProcedure
\end{algorithmic}
\end{algorithm}

}
